\quellseite{80}
\chapter{B. Religious Instruction in the Middle Ages}

J. W. Hoefling i, pp. 551 ff.; ii, pp. 27–49. *G. v. Zezschwitz i, ch. 19. 23; ii, 1, ch. 11. 18. 25. 27. 28. 30. 32. 33; ii, 2, 1, ch. 11; ii, 2, 2, ch. 30. Fr. Probst, \emph{Geschichte der katholischen Katechese}, pp. 55–129. Th. Harnack, pp. 55–65. G. Schumann and E. Sperber, \emph{Gesch. des Religionsunterrichts}, 1890, pp. 14–23. *E. Sachsse, pp. 116–185. F. Cohrs, \emph{Katechismusunterricht im Mittelalter} in PRE³, 1901. W. Caspari, \emph{Konfirmation} in PRE³, 1901. E. Chr. Achelis, pp. 294–303. G. Rietschel, \emph{Liturgik} ii, pp. 8–63. P. Drews, \emph{Taufe} in PRE³, 1907. J. Steinbeck, pp. 10–26. R. v. Raumer, \emph{Die Einwirkung des Christentums auf die althochdeutsche Sprache}, 1845. G. Th. Ditmar, \emph{Beitraege zur Geschichte des katechetischen Unterrichts in Deutschland}, 1848. P. Goebl, \emph{Geschichte der Katechese im Abendland vom Verfall des Katechumenats bis zum Ende des Mittelalters}, 1880. A. T. Drane, \emph{Christian Schools and Scholars from the Christian Era to the Council of Trent}, 1881. *H. J. Holtzmann, \emph{Die Katechese des Mittelalters} in Zeitschrift fuer prakt. Theologie, 1898. O. Dibelius, \emph{Das Vaterunser. Umrisse zu einer Geschichte des Gebets in der alten u. mittelalterl. Kirche}, 1903. *Fr. Wiegand, \emph{Das apostolische Symbol im Mittelalter}, 1904. K. A. Schmid, \emph{Geschichte der Erziehung} ii, 1, 1892. *A. Hauck, \emph{Kirchengeschichte Deutschlands} i ³1904, ii ³1912, iii ³1906, iv ²1913, v 1, 1911. F. Watson, \emph{English Grammar Schools Before 1660}, 1909. F. P. Graves, \emph{History of Education during the Middle Ages}, 1909. A. F. Leach, \emph{Church Schools, Bishop Schools, Cathedral Schools} in Monroe’s Cyclopedia of Education, 1911. A. F. Leach, \emph{Schools of Medieval England}, 1915.

\quellseite{81}
\section[§ 11. From the Seventh Century]{§ 11. Religious Instruction from the Seventh Century to the End of the Twelfth}

*A. Hauck, \emph{Kirchengeschichte Deutschlands} i, pp. 208–216. 341 ff. 368. 394. 492. 531 f. 552. 556. 569; ii, p. 123. 168 f. 185 ff. 237 ff. 270 ff. 392. 397. 466 f. 724; iii, pp. 323 ff. 929 ff.; iv, p.38. *Fr. Wiegand, \emph{Die Stellung des apostol. Symbols im Mittelalter} i, pp. 262–351. W. J. Townsend, \emph{The Great Schoolmen of the Middle Ages}, 1881. F. A. Specht, \emph{Geschichte des Unterrichtswesens in Deutschland von den aeltesten Zeiten bis zur Mitte des 13. Jahrhunderts}, 1885. J. Baier, \emph{Der heilige Bruno, Bischof von Wuerzburg, als Katechet}, 1893. Fr. Wiegand, \emph{Erzbischof Odilbert von Mailand ueber die Taufe}, 1899. A. F. West, \emph{Alcuin and the Rise of the Christian Schools}, 1892. H. Dietscheid, \emph{Alkuins Leben und Bedeutung fuer den religioesen Unterricht} ii, 1893. *M. Manitius, \emph{Geschichte der lateinischen Literatur des Mittelalters}, vol. i: Von Justinian bis zur Mitte des 10. Jahrhunderts, 1911. P. Sprockhoff, \emph{Althochdeutsche Katechetik}, 1912. G. H. Hoerle, \emph{Fruehmittelalterliche Moenchs- und Klerikerbildung in Italien}, 1914. \emph{Dicta abbatis Pirminii} (sic!) de singulis libris canonicis \emph{scarapsus}, reprinted in a reliable text in P. Caspari, \emph{Kirchenhistorische Anekdota}, 1883, pp. 151 ff. \emph{Bonifatii epistolae}, ed. by E. Duemmler, \emph{Monumenta Germaniae Historica}, \emph{Epistolae}, vol. iii, pp. 215 ff. \gesperrt{CHARLEMAGNE}: \emph{Capitularia regum Francorum}, MGH i, 1883; \emph{Concilia}, MGH ii, 1906; Epistolae, MGH iii and iv, 1892, 1895; *H. Schuetze, \emph{Auslese aus den Werken beruehmter Lehrer und Paedagogen des Mittelalters}, No. 5: *\emph{Verordnungen und Briefe von Karl d. Gr.} (Charlemagne's circular letter to the archbishops, p. 26; cf. also Wiegand, \emph{Erzbischof Odilbert} etc., pp. 23 ff.; on p. 9 Wiegand notes where the answers of the bishops can be found). \emph{Alcuini opera} in MSL c and ci; *\emph{Disputatio puerorum per interrogationes et responsiones} in vol. c, pp. 1099 f.; German in Probst, \emph{Geschichte der katholischen Katechese}, pp. 89 ff.; ch. 2 contains an explanation of the Creed, and ch. 12 an explanation of the Lord's Prayer, both explanations in German and Latin in Baier, \emph{Bruno v. Wuerzburg}, pp. 117 ff. Alcuin's letters also in MGH, Epistolae iv, ed. by E. Duemmler, 1895. H. Schuetze, Auslese etc., No. 3: \emph{Alkuins Brief an Karl den Grossen}, 1879. *J. Freundgen, \emph{Alkuins paedagogische Schriften}, 1899. Ratio de catechezandis rudibus and \emph{Baptismal Discourses} by Maxentius of
\quellseite{82}
Aquileia and by an anonymous writer are edited by J. M. Heer, \emph{Ein Karolingischer Missionskatechismus}, 1911. \emph{Rhabani Mauri opera} in MSL cvii-cxii; \emph{De institutione clericorum} cvii, pp. 295 ff.; *\emph{De ecclesiastica disciplina} (of special importance) cxii, pp. 191 ff.; De institutione clericorum in separate edition by A. Knoepfler, 1900. *\emph{J. Freundgen, Rhabanus Maurus’ paedagogische Schriften}, 1889. Jonas of Orleans; \emph{De institutione laicali} in MSL cvi, pp. 152 ff. (In i, 6.8 and ii, 14 he treats of the duties of sponsors and parents in regard to the Christian training of their children.) Dhuoda (or Dodana), \emph{Liber manualis} in MSL cvi, pp. 109 ff.; a more complete edition by Bondurant 1887; German by G. Meier in \emph{Bibliothek der katholischen Paedagogik} iii, 1890 (Dhuoda wrote this interesting instruction in proper Christian conduct for her son who at the time was absent from home; it was begun in 841). *K. Muellenhoff u. W. Scherer, \emph{Denkmaeler deutscher Poesie und Prosa aus dem 8.-12. Jahrhundert}, third edition by E. Steinmeyer, 2 vols. 1892: \emph{Exhortatio ad plebem christianam} i, pp. 200–201; \emph{Freising explanation of the Lord’s Prayer} i, pp. 202–203; \emph{The Weissenburg Catechism} i, pp. 204–209; \emph{Notker’s Catechism} i, 249–259; also a large number of Old High German formulae of faith, prayer, and confession. *\emph{Heliand}, ed. by P. Piper, 1897; German by K. Simrock, 81882. *\emph{Otfried’s Gospel Book}, ed. by P. Piper, 1878; German by G. Rapp, 1858. *\emph{Theoduli Ecloga}, rec. J. Osternacher, 1904; cf. J. Osternacher, \emph{Quos auctores Latinos et sacrorum Bibliorum locos Theodulus imitatus esse videatur}, 1907; K. Knoke in \emph{Haus und Schule}, Hannoversches Zeitblatt, 1874, pp. 124 ff.; K. Knoke in \emph{Theologische Studien und Kritiken}, 1881, pp. 777 f.; H. Vollmer, \emph{Zur Geschichte des biblischen Unterrichts} in Baumgarten’s Monatsschrift fuer die kirchl. Praxis, 1904, pp. 321 ff.; H. Vollmer, \emph{Beitraege zur Geschichte des bibl. Unterrichts, bes. in Deutschland} in Mitteilungen der Gesellschaft fuer deutsche Erziehungs- und Schulgeschichte, 1904, pp. 278 ff.; *G. L. Hamilton, \emph{Theodulus, a Medieval Textbook in Modern Philology}, 1909; M. Manitius, loc. cit. i, pp. 568–574. *\emph{Comestoris Historia scholastica} in MSL ciic, pp. 1050 ff., 1855. \gesperrt{BIBLIA HISTORICA}: E. Reusch, \emph{Die deutsche Historienbibel vor Erfindung der Buchdruckerkunst} in Beitraege zur theol. Wissenschaft von Reusch und Kunitz, 1855. *Th. Merzdorf, \emph{Die deutschen Historienbibeln des Mittelalters}, 2 vols., Texts, 1870; *H. Vollmer, \emph{Ober- und mitteldeutsche Historienbibeln}, 1912.—For \emph{Biblia Pauperum} compare literature for chapter 12. \gesperrt{MIRACLE AND MORALITY PLAYS}: F. J. Mone, \emph{Altdeutsche Schauspiele}, 1841; id., \emph{Schauspiele des Mittelalters} in \emph{Deutsche Nationalliteratur}
\quellseite{83}
xiv, 1891; Wilken \emph{Geschichte der geistlichen Spiele in Deutschland}, 1872; E. K. Chambers, \emph{The Medieval Stage}, 2 vols., 1903; W. Creizenach, \emph{Geistliche Spiele} in PRE³, 1906.

The Church of the Middle Ages found herself face to face with two problems the solution of which was to determine her entire subsequent development. The one problem arose from the gradual dissolution of the catechumenate for adults in which, as was seen, Christian saving truth had been all but replaced by liturgical ceremonies. The other problem arose from the fact that so far no method of instruction for those baptized in infancy had been worked out.

Let us begin by turning our attention to the Church’s attempts at solving the former of the two problems. Very largely it was not seriously attacked or even appreciated. Quite generally the missionary preachers thought that they had fulfilled their duty if they impressively proclaimed the gospel and immediately baptized those whom they had thus convinced of the vanity of idolatry and of the greater worth of Christianity. This was the practice of Augustine when he labored among the Anglo-Saxons (597/8), of Wilfrid and Willibrord among the Frisians (678; 690 ff.), even of Boniface among the Hessians and Frisians (719 ff.), of Ansgar among the Danes and Swedes (837 ff.), and of Otto of Bamberg among the Pomeranians (1124–1128). They admitted to baptism whoever applied for it; prebaptismal instruction did not exist. Moreover, if we read how Boniface expressly enjoined that the candidate should be given the formula of renunciation and the Creed in his native tongue, we must conclude that even this was not always done and that therefore often the candidate did not know what he renounced or professed. And baptismal instruction was, of course, altogether out of the question when Christianity was by force imposed upon whole nations (as the Saxons under Charlemagne, the Livonians by the Order of the Sword, 1200 ff., or the Prussians by the Teutonic Knights, 1230 ff.).

\quellseite{84}
Frequently post-baptismal care was also sadly lacking. The rite of \emph{confirmatio} which in some localities had been detached from baptism and was allowed to be administered only by the bishop, was seldom properly utilized for instruction; and most of the priests were not equal to the task of remedying the neglect through preaching and instruction after baptism. The belief in the magical effect of the sacraments everywhere manifested its destructive power. It is evident that under such circumstances only a defective and barren kind of Christian life could result. These defects were intensified by the unscrupulous manner in which the well meant counsel of Gregory the Great was observed. He had advised the missionaries to retain wherever possible the ancient pagan forms and sanctuaries and to give them a Christian interpretation; but the result was an incongruous combination of worship true and false, of magic formula and Christian ceremony; yes, magic efficacy was only too often ascribed to Christian rites. Such were the conditions among the Alemanni near the Lake of Constance and in Alsace, according to Pirmin’s testimony, or, as Boniface tells, in Bavaria. Centuries were required to retrieve this grievous error.

To be sure, it would have been impracticable simply to transplant the catechumenate of the Early Church into this new environment. The catechumenate was developed in the highly civilized Graeco-Roman world, in an age of pronounced individualism where no one joined the Church unless he was inwardly, personally convinced of the truth of Christianity. Now, however, the Church dealt with primitive peoples as yet possessing no literature, with nations and tribes of the strongest collectivism where the individual simply followed the decision of the group, either for or against Christianity. But this profound difference did not absolve the Church from the duty of baptismal discipline.

A few voices in the wilderness laid stress upon instruction.
\quellseite{85}
Above all, we name \emph{Pirmin}. Whenever adults were to be baptized among the “christianized” Alemanni (724–753), he first instructed them, using for this purpose the “catalogue of mortal sins” (\emph{cupiditas, gula, fornicatio, ira, tristitia, acedia, vana gloria, superbia}; cf. Gal. 5:19 ff. and 1 Cor. 6:9 f.) and the Apostolic Creed. Instruction, in his mind, was the principal thing; in distinction from Boniface who introduced the complete Roman baptismal ritual, Pirmin employed a very plain form of baptism consisting of renunciation, confession of the Creed, immersion, anointment, and robing in white. A little later, it was \emph{Alcuin} who insisted on thorough prebaptismal instruction when the Christianisation of the Avari was undertaken; and Charlemagne remembering the evil consequences of the forcible conversion of the Saxons, accepted his advice. Alcuin reverted to the practice of the Early Church, recommending to the missionaries Augustine’s book \emph{De catechizandis rudibus}, as the most appropriate manual. He demanded that the course of instruction should offer general information concerning immortality, judgment, life eternal and damnation, a brief explanation of the symbol and the Lord’s Prayer, and after baptism, an introduction into the commandments of the gospel. The commission appointed by Charlemagne adopted his counsel. It also decided that the period of instruction should last no longer that 40 days nor less than 7 days, an echo, as it were, of the competent institute of the Early Church. It is true, the time allotted seems entirely too short, and it is also likely that the instruction did not always attain to the high level of the extant missionary catechism from the diocese of Ratisbon (about 800); nevertheless, it is commendable that the Church once more seriously attempted to do justice to her duty and succeeded in a measure far surpassing the achievements of the two previous centuries. Alcuin’s example was followed by \emph{Rhabanus Maurus}. In his letter to Reginbald, a missionary among the Slavs (about
\quellseite{86}
845), he prescribed—again in close agreement with Augustine—a course of baptismal instruction consisting of the \emph{narratio}, i.e., the essentials of sacred and church history, the \emph{exhortatio} based on the hope of the resurrection, and of an explanation of the Creed and the Lord’s Prayer. How widely and how long these directives were effectively obeyed, is unknown.

Except in mission territory, \emph{infant baptism} was, of course, universally practiced. In England, for instance, as early as 691 all children of Christian parents had to be baptized within thirty days under penalty of law, in Northumberland within nine days; Charlemagne fixed the limit at one year. In Germany the traditional baptismal dates, Easter and Pentecost, enjoyed special favor for five or six centuries. This was due to the introduction of the Roman service, which emphasized such externals. The baptismal rite was also performed according to the Roman form although the so-called Gregorian Sacramentary which Hadrian I sent to Charlemagne was not adopted without change. Immersion remained the more usual mode of baptism, but after 1287 aspersion, or sprinkling, gained ever greater prevalence.

This carries us to the second problem. How did the Medieval Church adjust herself to the problem of \emph{instructing those who were baptized in their infancy}? Such training was urgently needed, far more so than in the Early Church because of the almost negligible number of Christian homes. Yet, no attempt was made at introducing some sort of discipline; more yet, the seriousness of the situation was at first not even appreciated. In Alemannia, for instance, the clergy simply performed the magic rite of baptism and paid no more attention whatsoever to the baptized children. Not so \emph{Pirmin}. He nobly endeavored to instruct the children after their baptism. As material he used the catalogue of mortal sins and the Creed (now divided into 12 articles because of the legend
\quellseite{87}
that each of the apostles contributed his share at the time of its composition). Boniface likewise was active in this direction. But the achievements of these men were rather fragmentary. No fixed and stable method was established until the time of \emph{Charlemagne}. It is true, he enjoined that the ritual of the scrutinies should be followed in infant baptism and that at least the final scrutiny should be conducted completely; but he also demanded that instruction be imparted in connection with this ceremony, not of course to the children themselves, but to their \emph{sponsors}. They were addressed in the introductory catechesis, they were instructed in the Creed and the Lord’s Prayer, they had to learn them by heart and recite them in the place of the children, and had to promise solemnly that they would teach them to their god-children. In a stringent decree, Charlemagne barred from sponsorship all persons who were unable to recite these parts, and he took drastic steps against the recalcitrant. In this way the state took control of a task which properly belonged to the Church: no wonder that henceforth baptismal instruction savored of legalism and compulsion. The parents likewise were directed to instruct their children. Moreover, the priests were required to conduct periodical meetings of parents, to teach them the Creed and the Lord’s Prayer, and to exhort them to retain these parts in memory. The institute of auricular confession allowed at least some sort of control in this respect; those parents and sponsors who had failed to do their duty were excluded from Holy Communion and from sponsorship, and at their occasional visitations the bishops were required to examine especially this phase of congregational life. Likewise an attempt was made at enabling the clergy to fulfil their duties and at controlling them: every parsonage library, it was decreed, must contain an authorized and practicable explanation of the Creed and the Lord’s Prayer, and the bishop must ascertain whether it has been
\quellseite{88}
put to faithful use. Also in their sermons the priests had an opportunity to offer instruction, in fact most of the preaching, such as it was, turned upon the material supplied in the catalogue of mortal sins, the Creed, and the Lord’s Prayer, and a sermon especially on the articles of the Creed had to be delivered on the last Sunday in lent in order that the baptismal rite performed on Easter Sunday might be intelligently appreciated.

The Creed and the Lord’s Prayer were translated into German for the benefit of the pastors and congregations and were made the basis of doctrinal and edifying tracts. Thus the \emph{first German catechisms} came into existence. The most important of them is the so-called Weissenburg Catechism of 789. It contains the following five parts: (1) The Lord’s Prayer in German with a brief explanation based on Tertullian and Cyprian (partly appropriated, by the way, by Luther); (2) an enumeration of mortal sins according to Gal. 5:19–21, in Latin and German; (3) the Apostolic Creed, German and Latin; (4) the Athanasian Creed, German and Latin; (5) the \emph{Gloria in Excelsis}, German and Latin. Other noteworthy catechisms dating from this period are the Freising Explanation of the Lord’s Prayer (802); an erotematic explanation of Creed and Lord’s Prayer attributed to Alcuin (\emph{Disputatio puerorum per interrogationes et responsiones}, ch. xi and xii); and the catechism of Notker of St. Gall (950–1022) which contains, in both German and Latin, the texts of the Lord’s Prayer and of the two Creeds, and German explanations of these parts (only the explanation of the Lord’s Prayer is of moment).

The seriousness of purpose and the energy of effort revealed in these labors are certainly commendable, especially when we bear in mind the tremendous difficulties connected with the task of instructing those baptized in infancy. But it must be admitted that rarely more than rote learning was
\quellseite{89}
achieved, and that the dangerous substitution of knowledge for real, saving faith was very effectively promoted. This manner of externalizing what lives in the innermost recesses of the heart manifested itself already in the Early Church; now, however, it had well nigh succeeded in crushing true inward piety.

Owing to Charlemagne’s ceaseless efforts somewhat better results were attained in the \emph{schools} which were extended and raised to a higher standard of efficiency. This does not apply so much to the “parish” schools; here only those boys received more training who were to become attendants at the services. But in the cathedral and convent schools the Lord’s Prayer, the Apostolic and the Athanasian Creeds were learned, Psalms were sung and memorized, and the regular New Testament pericopes were studied. The pupils were also made acquainted with some of the most important phases of sacred history: the explanation of the Creeds was somewhat historical in tenor, the sermons — since Charlemagne preaching became more general — often contained material from sacred history, and the prevalent exercise of calculating the dates of festivals gave frequent occasion for such study. Even more useful in this respect was the widely used \emph{Ecloga} written by Theodul (i.e., Gottschalk) about the middle of the ninth century; this was a dialogue in verse between a young man personifying the spirit of paganism and a young woman representing Christianity, and offered much historical material, chiefly from the Old Testament. Some knowledge of the life of Christ undoubtedly was communicated by the epic \emph{Heliand} (830) and the \emph{Krist} (860; Otfrid of Weissenburg is the author). But genuine instruction in Bible history was not given in these schools, and after all their sphere of influence was very limited in as much as they were chiefly attended by future clerics and \emph{religiosi}.

As soon as these institutions passed from under the
\quellseite{90}
vigorous hand of Charlemagne, many features fell into disuse in spite of the earnest efforts of Louis the Pious; the fundamental ones, however, remained in force and in subsequent times were now and then employed by synods and councils as the basis for further development. Numerous decrees show that the teaching duty was often impressed upon parents, sponsors, and priests. Alcuin’s explanation of the Creed and the Lord’s Prayer was in current use in Southern Germany for over two centuries, and Notker’s catechism was used until the twelfth century. More than 150 manuscript copies of Theodul’s \emph{Ecloga} are extant, the copies dating from every century of the Middle Ages, and we possess a printed copy dated 1504. In the latter part of the twelfth century, the \emph{Scholastica Historia} by Peter Comestor (†1179) was added to this class of literature; this bulky work, however, was not intended for the instruction of children and reached only a very limited number of people. Since the eleventh century some knowledge of Biblical history was also imparted by means of religious spectacles and by mural, altar, and window decorations in the churches (Note that the \emph{Biblia Pauperum} to be mentioned later, originated from the headpiece of the altar of Klosterneuburg, 1181). But when all has been said, the fact remains that baptized children grew up without receiving regular instruction in Bible History and many of them without any instruction at all. Moreover, what little religious information was imparted, was largely contaminated with apocryphal and legendary material of every description.

\quellseite{91}
\section[§ 12. Religious Instruction to the End of the Middle Ages]{§ 12. Religious Instruction from the Thirteenth Century to the End of the Middle Ages}

*J. Geffken, \emph{Der Bilderkatechismus des 15. Jahrhunderts und die katechetischen Hauptstuecke in dieser Zeit bis auf Luther}, vol. i, \emph{Die zehn Gebote}, 1855. *V. Hasak, \emph{Der christliche Glaube des deutschen Volkes beim Schluss des Mittelalters, dargestellt in deutschen Sprachdenkmalen}, 1868. H. Brueck, \emph{Der religioese Unterricht fuer Jugend und Volk in Deutschland in der 2. Haelfte des 15. Jahrhdts}, 1876. F. Hipler, \emph{Christliche Lehre und Erziehung im Ermeland und im preussischen Ordensstaat waehrend des Mittelalters}, 1877. Chr. Moufang, \emph{Die Mainzer Katechismen von Erfindung der Buchdruckerkunst bis zum Ende des 18. Jahrhunderts}, 1877. J. Janssen, \emph{Geschichte des deutschen Volkes seit Ausgang des Mittelalters} i, ¹1877. *C. A. H. Burckhardt, \emph{Geschichte der saechsischen Kirchen- und Schulvisitationen von 1524–1545}, 1879. Fr. Falk, \emph{Die Druckkunst im Dienst der Kirche, zunaechst in Deutschland, bis zum Jahr 1520}, 1879. F. Schneider, \emph{Die Tischgebete in Luthers Katechismus} in Anzeiger fuer Kunde der deutschen Vorzeit, 1880. K. Grube, \emph{Johannes Busch, ein kathol. Reformer des 15. Jahrh.}, 1881. *Chr. Moufang, \emph{Katholische Katechismen des 16. Jahrh.}, 1881. V. Hasak, \emph{Dr. M. Luther und die religioese Literatur seiner Zeit bis zum Jahre 1520}, 1881. H. Weber, \emph{Geschichte des Christenlehr-unterrichts und der Katechismen im Bistum Bamberg}, 1882. O. Kaemmel, \emph{Gesch. des deutschen Schulwesens im Uebergang vom Mittelalter zur Neuzeit}, 1882. *J. Mueller, \emph{Quellenschriften zur Gesch. des deutschsprachl. Unterrichts bis zur Mitte des 16. Jahrh.}, 1882. Th. Esser, \emph{Gesch. des englischen Grusses [Ave Maria]} in Histor. Jahrbuch der Goerres-Gesellschaft, vol. v, pp. 88–116, 1884. *J. Mueller, \emph{Schulordnungen und Schulvertraege in hoch- und niederdeutscher Sprache von 1296–1523}, 1885. I. H. Lupton, \emph{A Life of John Colet, with an appendix of some of his English writings}, 1887. *J. Mueller, \emph{Die deutschen Katechismen der boehmischen Brueder}, 1887. J. Mueller, \emph{Die Anfaenge des saechsischen Schulwesens} in N. Archiv fuer Saechsische Gesch. und Altertumskunde, 1887. *F. Falk, \emph{Der Unterricht d. Volkes in den katechetischen Hauptstuecken am Ende des Mittelalters} in Hist.-politische Blaetter, 1891, pp. 553 ff.; 682 ff.; 1892, 81 ff.; 721 ff.; 729 ff.
\quellseite{92}
*P. Bahlmann, \emph{Deutschlands katholische Katechismen bis zum Ende des 16. Jahr.}, 1894. A. Roesler, \emph{Kardinal Joh. Dominicis Erziehungslehre und die übrigen paedagogischen Leistungen Italiens im 15. Jahrh.}, 1894. J. Mueller, \emph{Zu den Schriften des Mag. Nicolaus Rutze in Rostock} in Zeitschr. fuer niedersaechs. Kirchengesch., 1896. F. A. Gasquet, \emph{Religious Instruction in England during the 14th and 15th Centuries} in \emph{The Old English Bible and other essays}, 1897. F. Cohrs, \emph{Zur Katechese am Ende des Mittelalters} in Zeitschr. fuer prakt. Theologie, 1898. *F. Cohrs, \emph{Die evangelischen Katechismusversuche vor Luthers Enchiridion}, vol. iv, 1902. J. Richter, \emph{Die paedagogische Literatur in Frankreich waehrend des 16. Jahrh.}, 1904. *F. Wiegand, \emph{Das apostolische Symbol im Mittelalter}, 1904. *F. H. Kunz, \emph{Joh. Gersons paedagogische Schriften} in Bibl. d. kathol. Paedagogik, vol. xv, 1904. F. H. Thalhofer, \emph{Die katechetischen Lehrstuecke im Mittelalter} in Mitteilungen der Gesellschaft fuer deutsche Erziehungs- und Schulgeschichte, 1905. W. Burger, \emph{Roemische Beitraege zur Gesch. der Katechese im Mittelalter} in Roemische Quartalschrift, 1907. *M. Reu, \emph{Quellen zur Geschichte des kirchl. Unterrichts im evgl. Deutschland zwischen 1530 und 1600}, part 2, 1906. G. S. Coulton, \emph{Mediaeval Studies}, No. 7: \emph{Religious Education before the Reformation}, 1906. *J. Greving, \emph{Zum vorreformatorischen Beichtunterricht}, 1907. *R. Galle, \emph{An der Wiege des biblischen Geschichtsunterrichts und Luthers Passionale} in Mitteilungen etc., 1907. F. Felinger, \emph{Das Kind in der altfranzoesischen Literatur}, 1908. M. Reu, \emph{Der kirchl. Unterricht am Vorabend der Reformation} in Kirchl. Zeitschr., 1910. A. Hauck, \emph{Kirchengesch. Deutschlands} v 1, pp. 340 ff.; 358 ff., 1911. M. Reu, \emph{Religious Instruction of the Young in the Century before the Reformation} in Luth. Church Review, 1915. \emph{Biblia Pauperum}: Camesina u. Heider, \emph{Biblia Pauperum}, Darstellungen der \emph{Biblia Pauperum} in einer Handschrift des 14. Jahrh., 1863. Laib u. Schwarz, \emph{Biblia Pauperum}. Nach dem Original zu Constanz, ²1892. W. L. Schreiber, \emph{Die Entstehung und Entwicklung der B. P. unter bes. Beruecksichtigung der uns erhaltenen Handschriften}, 29 Textillustrationen und 11 Lichtdrucktafeln, 1903. *P. Heitz u. W. L. Schreiber, \emph{Biblia Pauperum}. Nach dem einzigen Exemplar in 50 Darstellungen, 1903. F. J. Luttor, \emph{Biblia Pauperum}, Studie zur Herstellung eines inneren Systems. Mit dem Text der in der Wiener Hofbibliothek aufbewahrten Handschrift, 1912. J. Lutz and P. Perdrizet, \emph{Speculum humanae salvationis}, critical edition, 1906 f. *H. Vollmer, \emph{Eine deutsche Schulbibel des 15. Jahrhunderts. Historia scholastica des Petrus Comestor in deutschem Auszug mit latein.}
\quellseite{93}
Paralleltext, 1926. \emph{Abaelardi opera}, MSL clxxviii; \emph{Explanation of the Apostolic Symbol}, pp. 617–630. Thomas Aquinas, \emph{Opuscula}, Paris 1656; opusculum v: \emph{De articulis fidei et ecclesiae sacramentis}, pp. 47–54; op. iv: \emph{De decem praeceptis et lege amoris}, pp. 30–46; op. vii and viii: \emph{Expositio orationis dominicae et salutationis angelicae}, pp. 206–218; op. xvi: \emph{Expositio Symboli Apostolorum}, pp. 240–258. C. Martin, \emph{Ausgewaehlte Schriften des Thomas v. Aquin} (Two commandments of love, the ten commandments of God, and the holy sacrament of the altar), 1852. *A. Portmann und X. Kunz, \emph{Katechismus des hl. Thomas v. Aquin}, ²1899; here we find five other catechetical writings intended chiefly for adults or even for theologians: (1) \emph{Speculum ecclesiae} by Edmund of Canterbury († 1242), pp. 261–316; (2) \emph{La Somme le Roi} by Laurentius Gallus († 1285), pp. 317–337; (3) \emph{Summary of Christian Doctrine} according to the acts of the Synod of Toledo (1323), pp. 342–346; (4) \emph{Trattato di dottrina cristiana} (first part of 14th century), pp. 349–384; (5) \emph{Fundamentum aeternae felicitatis} (14th or 15th century), pp. 388–405. \emph{Summary of Doctrine of the Synod of Lavaux} (1368), German in Neues Hedwigsblatt, Breslau, 1875. Richard Hampole († 1349), Bibliotheca Maxima (Lugd.) vol. xxvi: \emph{Explanatio orationis dominicae et symboli}, pp. 619–624. \emph{Symbolum et omnium credendorum ad salutem spectantium compendiosa collectio} is printed in Gersonii opera i, pp. 233–422, though not his own, cf. Geffken, p. 30. Arnold of Villanova: \emph{Alphabeticum catholicorum} (13th century), reprinted in Burger, pp. 173–194. \emph{Tabula fidei christianae} (14th century), ibid., pp. 194–197. \emph{Joannis Gersonii opera omnia}, 5 vols., 1706. J. B. Schwab, \emph{Johann Gerson}, 1858. *Gerson, \emph{De parvulis ad Christum trahendis}, ed. by N. Puengel, 1853; German by H. Schuetze, \emph{Auslese} (see literature for ch. 11), No. 1, 1879; another German translation with a good introduction by F. X. Kunz in Bibliothek der katholischen Paedagogik, vol. xv, 1904. Gerson, \emph{Opusculum tripartitum de praeceptis decalogi, de confessione et de arte moriendi} in opera i, pp. 425–450; German edition by Geiler of Kaiserberg under the title: \emph{Der dreieckicht Spiegel}, Strasburg, 1510 and 1514, edited in modernized form by V. Hasak in \emph{Herbstblumen oder alte ernste Wahrheiten. Zur Illustration des christlichen Volksunterrichts in der vorreformatorischen Zeit}, 1885, pp. 89–131. Gerson, \emph{A B C des simples gens} (a real catechism for children!), printed for the first time from a Paris manuscript by H. Jadart in Jean de Gerson, 1881, pp. 252–253. \emph{Symbolum Apostolicum oder das apostolische Glaubensbekenntnis mit beigefuegtem Text der Glaubensartikel}. 12 Holzschnitte in Faksimile-Reproduktion nach dem
\quellseite{94}
einzigen bis jetzt bekannten Exemplar der Hof- und Staatsbibliothek zu Muenchen, Stuttgart, s. a.; belongs to the 15th century; cf. the Heidelberg illustrations for the commandments, confession, and the seven mortal sins in Geffken, op. cit., suppl. Tafel des christlichen Lebens (from the end of the 15th century), reprinted by Bahlmann, pp. 61–75. John Colet: \emph{Catechison}, reprinted in J. H. Lupton, \emph{A Life of John Colet}, 1877, pp. 286 ff., also in Cohrs iv, pp. 421 ff., and by Reu in Kirch. Zeitschrift, 1910. Desiderius Erasmus: \emph{Institutum hominis Christiani}, reprinted in Cohrs iv, pp. 421 ff. Tritonius: \emph{Enchiridion}, cf. Cohrs iv, pp. 430 ff. \emph{Rostocker Kinderlehre} (first decades of the 15th century) reprinted by K. E. H. Krause, \emph{Programm der grossen Stadtschule zu Rostock}, 1873. Dietrich Koelde: \emph{Spegel des cristene mynschen} (1470), translated into High German by Moufang, \emph{Kathol. Katechismen}, pp. 1–1. On other popular writings containing catechetical explanations, e.\,g., \emph{Der Seele Trost}, Stephanus Lanzkranna’s \emph{Die Hymmelstrass}, \emph{Das Buch der 10 Gebote}, \emph{Spegel der leyen}, see Bahlmann and especially Geffken. *Ulrich Surgant: \emph{Manuale Curatorum}, Basel, 1503; esp. Libri secundi consideratio quinta (print of 1506, leaves 80–81) prescribes the form in which the chief parts are to be read from the pulpit every Sunday, reprinted in part by Reu, \emph{Religious Instruction} etc. On the Tabulae Catechismi still extant compare Reu, \emph{Der Religionsunterricht am Vorabend usw.}; also a Cologne print of about 1490 with the title \emph{Paternoster, Ave Maria, Glaubensbekenntnis und die 10 Gebote}; another single leaf print (Low German) in J. Halle, Antiquariat, Muenchen, \emph{Katalog xlvi}, No. 463 [Price \$125.00]. CONFESSIONAL BOOKLETS: Muenzenberger, \emph{Das Frankfurter und Magdeburger Beichtbuechlein und das Buch vom sterbenden Menschen}, 1881. *F. W. Battenberg, \emph{Das Beichtbuechlein des Magisters Joh. Wolff}, 1907. *F. Falk, \emph{Drei Beichtbuechlein nach den 10 Geboten aus der Fruehzeit der Buchdruckerkunst}, 1907. Compare also the rich material given in Geffken, op. cit., suppl. pp. 1–218. — Wiclif: \emph{His Works}, ed. by the Wiclif Society, 1883 ff.; \emph{Opera minora}, 1913, No. 5: \emph{De fide catholica}; No. 15: \emph{De oratione dominica}; No. 16: \emph{De salutatione angelica}; cf. Loserth, \emph{Wiclif in PRE³}, 1908. \emph{Writings of J. Wiclif}, ed. by Religious Tract Society, London, 1831; here is found the \emph{Pauper rusticus, or The Poor Caitiff}, i.\,e., a popular explanation of the customary parts of the catechism, viz., \emph{On the belief}, pp. 51–59; \emph{On the Ten Commandments}, pp. 60–86; \emph{On the Lord’s Prayer}, pp. 86–98; cf. Geffken, op. cit. suppl. pp. 214 ff. \emph{John Hus}: Since 1913 his works are being edited by Flasjhans. Of these the first volume
\quellseite{95}
contains the \emph{Expositio Decalogi}. Two works written in Bohemian are of special importance: (1) \emph{The three stranded cord of faith, love, hope, to which every Christian should cling if he would gain eternal salvation}, 1412 (in the edition of his works in Bohemian by Erben it is found in vol. iii, pp. 152 ff.); (2) \emph{Explanation of the creed, the ten divine commandments and the Lord’s Prayer} (in Erben’s edition, vol. i, pp. 1–386). Besides his postil, these two are the most important works of Hus in Bohemian. Like all his works, however, these are based upon similar writings of Wiclif. Both works were translated into Low German and published by Mag. Nicholaus Rutze (not Rus) of Rostock († about 1508). The first of these two works of Hus was published according to Rutze’s text in 1886 by K. Nerger in \emph{Osterprogramm des Rostocker Gymnasiums}. J. Mueller proved that Rutze merely translated from Hus; cf. also Geffken, suppl., pp. 159 ff. Moravian Catechisms: *Jos. Mueller, \emph{Die deutschen Katechismen der boehmischen Brueder} in Monumenta Germaniae Paedagogica vol. iv. Waldensian Catechism, reprinted by G. v. Zezschwitz, \emph{Die Katechismen der Waldenser und Boehmischen Brueder}, 1863; cf. K. Mueller, \emph{Die Waldenser}, 1886.

A new era began to dawn in Germany and the rest of the Western world as a result of the Crusades and the splendid achievements of the Hohenstaufen dynasty (esp. Frederick II, 1215–1250). True, the Church came out victoriously from the conflict with the State; but she also realized the imperative necessity of satisfying the growing desire for culture, of adjusting religious instruction to the demands of the new age. Besides, it was advisable no longer to ignore the teaching duty, for heretical movements arose and spread at an alarming rate. Charlemagne’s institutions were drawn forth from oblivion and sincere efforts were put forth to introduce them anew.

Much attention was paid (1) to the restoration of the domestic catechumenate. Parents and sponsors were strongly urged to instruct their children. Theodulf of Orleans in De ordine baptismi (about 788) had expressed the idea that instruction should be given when the “period of discernment” is reached; now the seventh year was fixed as the time when
\quellseite{96}
the young should be introduced to religious truth. In one of the sermons of Berthold of Ratisbon († 1272) we read: “When the child is seven years old, his sponsors shall teach him the Creed and the Lord’s Prayer; that is their duty, they are his spiritual parents. They shall say to his parents: ‘Friends, you must teach my god-child the Lord’s Prayer and the Creed; else he shall come to me and I will teach him.’ And if he learns also the Ave Maria, that is so much better.” Neglect of this duty is often found among the sins enumerated in the customary formulae of confession used in that period. The 13th century and still more the 14th have produced an abundance of ascetical tracts and books which helped to revive the domestic catechumenate. The more important are \emph{La Somme le Roi} by Laurentius Gallus (1279), \emph{Eyn speyghel des christen gelouben} by Ludolf of Goettingen (1472), \emph{Der Seele Trost}, author unknown (before 1472), \emph{Die Hymelstrass} by Stephan Lanzkranna (1484), \emph{Di dottrina cristiana} (13th century), \emph{Instructions of the Synod of Toledo} (1323), and especially John Gerson’s \emph{Opusculum tripartitum de praeceptis decalogi, de confessione et de arte moriendi} (about 1412). The last three works are designed for children; all of them were widely circulated, frequently copied and reprinted, and translated into various languages.

(2) During the preceding period confirmation had been separated from baptism (see p. 62); it was performed only by bishops, and gradually gained sacramental acceptation (probably Hugh of St. Victor, † 1144, was the first to count it as the second sacrament; by 1439 it was universally recognized). The age at which this rite was usually performed was the seventh year (so decided by the synod of Cologne, 1279, and the synod of Liege, 1287), but we have no information to what extent, if at all, instruction was connected with it. The confessional, however, offered an opportunity to supplement domestic instruction: the Fourth Lateran Council
\quellseite{97}
(1215) required all children from the seventh year to confess at least once a year and to commune on Easter Sunday (the communion of infants had gradually fallen into disuse in the Western Church, see p. 53); thus the priests were enabled to determine the amount, and the kind of instruction which the children had received at home, and often enough they themselves were impressed with the duty of supplementing the children’s information. The technical term applied to these confessional questions concerning Creed, Lord’s Prayer, etc., was \emph{catechismus}; we observe that this word still retained its original significance of oral teaching (cf § 1), but that it acquired the connotation of erotematic instruction, that is to say, instruction in the form of questions and answers. In England, France, and Germany, councils now and then also decreed that the clergy should teach the young the Creed and the Lord’s Prayer. The Ave Maria is mentioned in this connection with increasing frequency since mariolatry became popular, i.\,e., since the 13th century. The parents were exhorted to bring their children to church regularly every Sunday for the purpose of instruction. In order to equip the clergy for this work of teaching, special manuals, popular as well as technical, were prepared. We mention \emph{Expositio symboli apostolorum} by Thomas Aquinas (1274), \emph{Speculum ecclesiae} by Edmund of Canterbury (1240), \emph{Fundamentum aeternae felicitatis} by an unknown author (about 1470), and \emph{Beichtbuechlein} by John Wolff (1468). In certain localities, e.\,g., at Torgau, Basle, and Bamberg, the priests were even commanded to read, from the pulpit or from the “ambo”, the principal parts of instruction every Sunday. Ulrich Surgant’s widely circulated \emph{Manuale curatorum} (Basle, 1503) contains the form in which these parts were read to the congregation. Frequently the Creed, the Lord’s Prayer, the Decalogue, and other parts were written or printed on paper or cardboard and hung on the walls of churches, schools, or
\quellseite{98}
hospitals. Knowledge of this material was spread also by “catechetical sermons” which we meet occasionally toward the end of the Middle Ages.

(3) Religious instruction was imparted also in the \emph{Latin Schools} which were established in cities since the 14th century, and in the convent schools. In the lower classes the pupils had to memorize the Lord’s Prayer, Creed, and Ave Maria, now and then also the Benedictus and Gratias and the Confiteor, more rarely the Ten Commandments and a few liturgical formulae and hymns; in the upper classes some of these parts were learned in Greek. The same is true of the \emph{German schools for reading and writing} which since the 15th century were established throughout Germany, even in villages, under the influence of the Hanse and of the Brethren of the Common Life, except that here naturally these subjects were taught in the German language. One of the summaries of instruction for Latin schools has come down to our times, viz., the \emph{Catechyzon} by John Colet of London (1510). We also name two works which were intended especially for the instruction of children in the home (perhaps also for school use?), viz., \emph{A B C des simples gens} by John Gerson († 1429) and the Low-German \emph{Tafel des christlichen Lebens} by an unknown author (latter part of the 15th century).

While the catechetical parts received considerable attention in the schools, \emph{Bible History} was neglected, comparatively speaking. Information of only the most general type was imparted by means of mural and window paintings, by the observance of the church year, or by sermons on gospel texts. More was achieved by the widely spread \emph{Passionalia} and the \emph{Biblia Pauperum} (Old Testament and New Testament pictures placed side by side with brief explanations), and by Biblical picture books such as the \emph{Speculum humanae salvationis}. The most essential facts of sacred history were thus made known to old and young.

\quellseite{99}
If, however, we conclude from the above enumeration of facts that the Medieval Church did justice to her task of teaching or “ushered in a golden period of education,” we should commit a grievous error. True, the domestic catechumenate existed by right, but it was woefully deficient. The Church Visitations in Saxony during 1528 and 1529 disclose the sad fact that large numbers of adults knew neither Decalogue nor Creed nor Lord’s Prayer; what kind of instruction would such parents impart to their children? True, the parts of the catechism were read at some localities by the priest after the sermon; but that the rule was more frequently violated than observed appears from the prevalence of synodical decrees repeating it. Also the recorded minutes of visitations contain abundant and unambiguous evidence of the wide-spread lack of such reading. Besides, this reading was intended primarily for the benefit of adults, not of children. True, the priest was expected, by means of auricular confession, to exercise control over domestic instruction, and if necessary to supplement it; but how largely this order must have passed unnoticed appears when we read, e.\,g., in the official records of church visitations that a priest applied for a position who did not know the Ten Commandments and since six years had not looked into a book. True, in the Latin schools the catechism was incorporated into the text books; but this was done chiefly in order to preserve a connecting link between the German home and the Latin school, and not for the purpose of religious instruction. Moreover, a few notable cases excepted, the learning process in these schools consisted in mechanical memorizing. True, numerous expositions of catechetical material and Biblical picture books had been published; but their cost was so high and the number of illiterates so large that they affected only very limited circles. Besides, in point of contents, this whole literature was by no means evangelical—a few genuine
\quellseite{100}
gems excepted, but largely legendary, apocryphal, and even antibiblical.

Memory material was monstrously increased: it made hearing confession a much easier task for the priest. Besides the Lord’s Prayer, the Creed, the catalogue of mortal sins, and the Ave Maria, the young had to memorize the seven charisms, the seven cardinal virtues, the seven sacraments, the seven works of mercy, the eight Beatitudes, the twelve fruits of the Spirit, the Ten Commandments of God, the crying sins, the alien sins, the five senses etc. As the material increased, the possibility of inner appropriation decreased; the only way to arrive at mastery over the wealth of subject matter was mechanical memorizing with the aid of mnemonic verses or rhymes. The one gratifying feature is the final incorporation of the Decalogue in the catechism; henceforth it enjoyed, especially owing to its connection with confession, great prominence. As G. v. Zezschwitz said: “John the Baptist’s penitential preaching preceded the gospel preaching of the Reformation.”

By linking instruction with the confessional institute the Church completely shifted her educational aim. Confession being compulsory throughout life, the goal no longer was maturity and independence of the individual catechumen, but absolute submission to the clergy, lifelong subjection to these mediators between God and man.

Such was the result of the Church’s educational endeavors in the later Middle Ages. Heretical movements, such as the Waldensian, Wiclifite, or Hussite, with their emphasis of instruction and training, constituted a vigorous appeal to the Church to fulfill her sacred duty—but the call was not heeded. The earlier Waldenses (1170 ff.) taught their children the Decalogue and the Lord’s Prayer, also brief statements concerning the Trinity, the Church, the Sacraments, and the works of mercy, and insisted especially on cultivating Bible
\quellseite{101}
knowledge (e.\,g., they had the young memorize the current gospel and epistle lessons). \emph{Wiclif} did not write anything of a juvenile character; but his tracts on the Creed, Decalogue, and Lord’s Prayer (found in the \emph{Pauper rusticus, poor caitiff}) which were written for the home undoubtedly exerted some influence upon the domestic catechumenate. Nor did \emph{Hus} ever write books for the instruction of children; but two of his works (both of them dependent on \emph{Wiclif}) were used in the composition of the first Hussite Catechism (between 1420 and 1436) and the so-called Raudnitz Catechism (contemporary with the former). These two treatises by \emph{Hus} are entitled \emph{Das dreigeflochtene Tau aus Glaube, Liebe, Hoffnung} (=The Three-stranded cord of faith, love, hope), and \emph{Expositio symboli, decalogi et dominicae orationis}. The Hussite and the Raudnitz Catechisms form the basis for the first Catechism of the \emph{Moravian Brethren} (Bohemian before 1502, German 1522) entitled \emph{Ein christliche Unterweisung und Vorschrift der Jungen im Glauben} or simply \emph{Die Kinderfragen}, and also for the Catechism of the \emph{Waldenses} in Bohemia which bears the title \emph{Las interrogacions menors} (=The minor questions). Both of these catechisms bear the mark of a threefold division: faith, love, hope; they contain Symbol, Decalogue, Beatitudes, Lord’s Prayer, Sacraments, and other parts, and lay stress upon living, personal faith. The actual value of these catechisms is not so great; the \emph{Waldenses} and Moravians excelled the official Church rather by the faithful exercise of the parental teaching obligation. In this respect the Church would have done well to emulate their example.
